\documentclass[10pt, aspectratio=169]{beamer}
\usepackage{../beamer_theme_408}
\title{408考研零基础 --- 计算机网络}
\subtitle{4-12 传输层概述}
\author{}
\date{}
\begin{document}

\begin{frame}{本节目标}
    \begin{enumerate}
        \item 理解传输层的功能：端到端通信、进程到进程
        \item 理解端口号的概念与分类
        \item 掌握UDP协议特点及首部格式
        \item 掌握TCP协议特点，对比UDP
    \end{enumerate}
\end{frame}

\begin{frame}{传输层功能}
    \begin{block}{核心职责}
        提供\textbf{端到端}的通信服务（进程到进程），是应用层与网络层之间的桥梁。
    \end{block}
    \vspace{0.3em}
    \begin{center}
        \begin{tikzpicture}[>=stealth]
            \node[draw,minimum width=2.5cm,minimum height=0.6cm,fill=blue!20] at (-2,1.5) {应用层（进程）};
            \node[draw,minimum width=2.5cm,minimum height=0.6cm,fill=blue!20] at (2,1.5) {应用层（进程）};
            \node[draw,minimum width=2.5cm,minimum height=0.6cm,fill=red!20] at (-2,0.5) {\textbf{传输层}};
            \node[draw,minimum width=2.5cm,minimum height=0.6cm,fill=red!20] at (2,0.5) {\textbf{传输层}};
            \node[draw,minimum width=2.5cm,minimum height=0.6cm,fill=green!20] at (-2,-0.5) {网络层};
            \node[draw,minimum width=2.5cm,minimum height=0.6cm,fill=green!20] at (2,-0.5) {网络层};
            \draw[<->,thick] (-2,0.5) -- node[above]{端到端连接}(2,0.5);
            \draw[<->] (-2,-0.5) -- node[below]{逐跳转发}(2,-0.5);
            \draw[<->] (-2,1.2) -- (-2,0.8);
            \draw[<->] (2,1.2) -- (2,0.8);
        \end{tikzpicture}
    \end{center}
    \vspace{0.3em}
    \begin{columns}[t]
        \column{0.48\textwidth}
        \textbf{传输层提供的服务}
        \begin{itemize}
            \item 进程到进程通信
            \item 复用与分用
            \item 差错检测
            \item 流量控制
            \item 可靠传输（TCP）
        \end{itemize}
        \column{0.48\textwidth}
        \textbf{用端口号标识进程}
        \begin{itemize}
            \item IP地址标识主机
            \item 端口号标识进程
            \item 套接字 = IP + 端口
        \end{itemize}
    \end{columns}
\end{frame}

\begin{frame}{端口号}
    \begin{block}{端口号范围}
        16位端口号，范围0 -- 65535。
    \end{block}
    \vspace{0.3em}
    \begin{columns}[t]
        \column{0.33\textwidth}
        \textbf{熟知端口(0-1023)}
        \begin{itemize}
            \item HTTP: 80
            \item FTP: 21/20
            \item DNS: 53
            \item SMTP: 25
            \item SSH: 22
            \item HTTPS: 443
        \end{itemize}
        \column{0.33\textwidth}
        \textbf{登记端口(1024-49151)}
        \begin{itemize}
            \item 需IANA登记
            \item 避免冲突
            \item 例：8080(HTTP代理)
        \end{itemize}
        \column{0.33\textwidth}
        \textbf{动态端口(49152-65535)}
        \begin{itemize}
            \item 客户端临时使用
            \item 操作系统分配
            \item 通信完成后释放
        \end{itemize}
    \end{columns}
    \vspace{0.5em}
    \begin{center}
        \begin{tikzpicture}[>=stealth]
            \draw[->] (0,0) -- (8,0) node[right]{端口号};
            \draw[red,thick] (0,0.1) -- (1.5,0.1);
            \draw[blue,thick] (1.5,0.1) -- (6,0.1);
            \draw[green,thick] (6,0.1) -- (8,0.1);
            \node[align=center,above] at (0.75,0.4) {\small 熟知\\0-1023};
            \node[align=center,above] at (3.75,0.4) {\small 登记\\1024-49151};
            \node[align=center,above] at (7,0.4) {\small 动态\\49152-65535};
        \end{tikzpicture}
    \end{center}
\end{frame}

\begin{frame}{复用与分用}
    \begin{columns}[t]
        \column{0.5\textwidth}
        \textbf{复用}
        \begin{itemize}
            \item 多个应用层进程（如浏览器、邮件、FTP）
            \item 共用同一个传输层协议
            \item 封装时加上各自端口号
        \end{itemize}
        \vspace{0.5em}
        \begin{center}
            \begin{tikzpicture}[>=stealth]
                \node[draw,minimum width=1.5cm] (p1) at (-1.2,1.5) {进程1};
                \node[draw,minimum width=1.5cm] (p2) at (0,1.5) {进程2};
                \node[draw,minimum width=1.5cm] (p3) at (1.2,1.5) {进程3};
                \node[draw,minimum width=4cm] (t) at (0,0.5) {传输层};
                \draw[->] (p1) -- (t);
                \draw[->] (p2) -- (t);
                \draw[->] (p3) -- (t);
                \draw[->] (t) -- (0,-0.5);
                \node at (0,-1) {复用};
            \end{tikzpicture}
        \end{center}
        \column{0.5\textwidth}
        \textbf{分用}
        \begin{itemize}
            \item 传输层收到数据后
            \item 根据目的端口号
            \item 分发到对应应用进程
        \end{itemize}
        \vspace{0.5em}
        \begin{center}
            \begin{tikzpicture}[>=stealth]
                \node[draw,minimum width=4cm] (t) at (0,0.5) {传输层};
                \node[draw,minimum width=1.5cm] (p1) at (-1.2,-0.5) {进程1};
                \node[draw,minimum width=1.5cm] (p2) at (0,-0.5) {进程2};
                \node[draw,minimum width=1.5cm] (p3) at (1.2,-0.5) {进程3};
                \draw[->] (t) -- (p1);
                \draw[->] (t) -- (p2);
                \draw[->] (t) -- (p3);
                \draw[->] (0,1.2) -- (t);
                \node at (0,-1.2) {分用};
            \end{tikzpicture}
        \end{center}
    \end{columns}
\end{frame}

\begin{frame}{UDP协议特点}
    \begin{block}{UDP (User Datagram Protocol)}
        用户数据报协议，提供\textbf{无连接、不可靠}的传输服务。
    \end{block}
    \vspace{0.3em}
    \begin{columns}[t]
        \column{0.5\textwidth}
        \textbf{特点}
        \begin{itemize}
            \item \textbf{无连接}：发送前不建立连接
            \item \textbf{不可靠}：不保证送达，无确认重传
            \item \textbf{面向报文}：保留报文边界
            \item \textbf{开销小}：首部仅8字节
            \item \textbf{支持广播/组播}
        \end{itemize}
        \column{0.5\textwidth}
        \textbf{适用场景}
        \begin{itemize}
            \item DNS查询
            \item 视频/音频流媒体
            \item 实时游戏
            \item DHCP
            \item SNMP
        \end{itemize}
    \end{columns}
\end{frame}

\begin{frame}{UDP首部格式}
    \begin{center}
        \begin{tikzpicture}[>=stealth]
            \draw[thick] (0,0) -- (8,0);
            \draw[thick] (0,0) -- (0,-3.5);
            \draw[thick] (8,0) -- (8,-3.5);
            \draw[thick] (0,-3.5) -- (8,-3.5);
            \draw[thick] (4,0) -- (4,-0.8);
            \draw[thick] (0,-0.8) -- (8,-0.8);
            \draw[thick] (4,-0.8) -- (4,-1.6);
            \draw[thick] (0,-1.6) -- (8,-1.6);
            \draw[thick] (4,-1.6) -- (4,-2.4);
            \draw[thick] (0,-2.4) -- (8,-2.4);
            \node at (2,-0.4) {源端口(16位)};
            \node at (6,-0.4) {目的端口(16位)};
            \node at (2,-1.2) {UDP长度(16位)};
            \node at (6,-1.2) {UDP校验和(16位)};
            \node[draw,minimum width=8cm,minimum height=0.8cm,fill=gray!20] at (4,-2.8) {数据部分};
            \node[above] at (4,0.3) {\textbf{UDP首部（8字节）}};
        \end{tikzpicture}
    \end{center}
    \vspace{0.3em}
    \begin{itemize}
        \item 源端口：发送方端口（可选，不需应答时填0）
        \item 目的端口：接收方端口
        \item UDP长度：首部+数据的总长度（字节）
        \item 校验和：检测传输差错（可选，IPv4中可置0）
    \end{itemize}
\end{frame}

\begin{frame}{TCP协议特点}
    \begin{block}{TCP (Transmission Control Protocol)}
        传输控制协议，提供\textbf{面向连接、可靠}的传输服务。
    \end{block}
    \vspace{0.3em}
    \begin{columns}[t]
        \column{0.5\textwidth}
        \textbf{特点}
        \begin{itemize}
            \item \textbf{面向连接}：三次握手建立，四次挥手释放
            \item \textbf{可靠传输}：确认+重传
            \item \textbf{面向字节流}：不保留报文边界
            \item \textbf{全双工通信}
            \item \textbf{流量控制}（滑动窗口）
            \item \textbf{拥塞控制}
        \end{itemize}
        \column{0.5\textwidth}
        \textbf{适用场景}
        \begin{itemize}
            \item 网页浏览（HTTP）
            \item 文件传输（FTP）
            \item 电子邮件（SMTP）
            \item 远程登录（SSH）
            \item 数据库连接
        \end{itemize}
    \end{columns}
\end{frame}

\begin{frame}{TCP vs UDP对比}
    \begin{table}
        \centering
        \begin{tabular}{|c|c|c|}
            \hline
            \textbf{特性} & \textbf{TCP} & \textbf{UDP} \\
            \hline
            连接 & 面向连接 & 无连接 \\
            可靠性 & 可靠（确认重传） & 不可靠（尽力而为） \\
            数据边界 & 字节流 & 报文边界 \\
            首部大小 & 20字节 & 8字节 \\
            速度 & 慢 & 快 \\
            全双工 & 是 & 半双工（单工） \\
            流量/拥塞控制 & 支持 & 不支持 \\
            广播/组播 & 不支持 & 支持 \\
            典型应用 & HTTP/FTP/SMTP & DNS/视频/DHCP \\
            \hline
        \end{tabular}
    \end{table}
    \vspace{0.5em}
    \begin{exampleblock}{选择原则}
        \begin{itemize}
            \item 需要可靠传输 $\rightarrow$ TCP
            \item 实时性要求高、允许丢包 $\rightarrow$ UDP
        \end{itemize}
    \end{exampleblock}
\end{frame}

\begin{frame}{本节总结}
    \begin{enumerate}
        \item 传输层功能：端到端、进程到进程（端口号标识），复用与分用
        \item 端口号：熟知0-1023、登记1024-49151、动态49152-65535
        \item UDP：无连接、不可靠、面向报文、首部8字节
        \item TCP：面向连接、可靠、面向字节流、全双工、首部20字节
    \end{enumerate}
\end{frame}
\end{document}
